\subsection{Hľadanie DNA reťazcov}
% dorado -> clover ->Najvacsi cluster ->multialignment -> konsenzus primeru -> vyhladanie local alignmenty primeru v kazdom reade-> vybranie stringov za primerom -> volanie clover -> clustrov som zavolal multialignment -> konsenzus z toho clusteru ->(filtrovanie)-> mapa katic na ich vystupny signal  
%(Bez Omacky)
Hľadanie DNA reťazcov v nanopórových dátach je komplexný proces, ktorý zahŕňa niekoľko krokov. Najprv sa surové dáta získané z nanopórového sekvenovania prekladajú pomocou basecallera Dorado do formátu BAM, ktorý umožňuje ďalšie spracovanie. Následne sa použije klastrovací algoritmus Clover na zoskupenie sekvencií do klastrov na základe ich podobnosti na zaciatku retazca.
% Najvacsi cluster je primer ktory zacina kazdy retazec v readoch
Zo všetkých vytvorených klastrov sa vyberieme ten najväčší, nan zavolam moj naivny multialignment algoritmus, ktorý vytvorí konsenzuálnu sekvenciu pre tento cluster. Tento konsenzus slúži ako náš "primer", ktorý sa bude používať na vyhľadávanie v jednotlivých readoch.
% pre kazdy read() hladame lokalne zarovnanie s primerom a vyberiem sequenciu atomov za primerom
Pre každý read sa potom vyhľadá lokalné zarovnanie s primerom a vyberie sa sekvencia, ktorá nasleduje za primerom. 
Tieto sekvencie sa následne zoskupia pomocou Cloveru, čím vzniknú nové klastre.
Na každý z týchto klastrov sa opäť aplikuje multialignment, A zarovnané sekvencie sa použijú na vytvorenie konsenzuálnej sekvencie pre každý cluster.
tak ziskame dlhe retazce a z nich mozeme filtrovat tie ktore su dostatocne dlhe a dostacne zoskupene.

%Nakoniec sa vytvorí mapa katic (k-merov) na ich výstupný signál, čo nám umožní analyzovať pravdepodobnosti rôznych typov chýb pre každú bázu v reálnom reťazci. Táto mapa je kľúčová pre simuláciu chýb pri generovaní nových DNA sekvencií.